% Thuật ngữ chuyên ngành môn Các hệ thống CSDL (Glossary)
\newglossaryentry{csdl}
{
    name=Cơ sở dữ liệu,
    description={Tập hợp các dữ liệu có liên quan logic với nhau và được thiết kế nhằm đáp ứng nhu cầu khai thác thông tin của một tổ chức}
}

\newglossaryentry{csdlpt}
{
    name=CSDL Phân tán,
    description={Tập hợp nhiều cơ sở dữ liệu có quan hệ logic với nhau được phân tán trên các vị trí mạng máy tính khác nhau}
}

\newglossaryentry{bcnf}
{
    name=BCNF,
    description={Dạng chuẩn Boyce-Codd trong thiết kế CSDL quan hệ, yêu cầu mọi vế trái của phụ thuộc hàm không tầm thường đều phải là siêu khóa}
}

\newglossaryentry{trigger}
{
    name=Trigger,
    description={Thủ tục đặc biệt được hệ quản trị CSDL tự động kích hoạt thực thi khi xảy ra các thao tác INSERT, UPDATE, hoặc DELETE trên bảng}
}

\newglossaryentry{transaction}
{
    name=Transaction,
    description={Giao dịch, một chuỗi các thao tác CSDL được thực thi như một đơn vị công việc logic duy nhất thỏa mãn trọn vẹn tính chất ACID}
}
